\subsection*{NT-A04. General Hecke angle discrepancy and its norm-twist exception}
Let $K$ be a number field and let $\eta$ be a unitary Hecke character
which is \emph{not} a finite-order character times a unitary power
of the norm. For every prime ideal unramified for $\eta$, define
$\theta_{\mathfrak p}\in\mathbb R/\mathbb Z$ by
$\eta(\mathfrak p)=e^{2\pi i\theta_{\mathfrak p}}$.
Assume RH for $L_K(s,\eta^m)$ for all positive integers $m$, and RH
for $\zeta_K$. Then, uniformly over all circle arcs $J$,
\[
 \#\{\mathfrak p\text{ unramified}:N\mathfrak p\le x,\
                          \theta_{\mathfrak p}\in J\}
 =|J|\operatorname{Li}_2(x)
       +O_{K,\eta}\bigl(\sqrt x\log^2(2x)\bigr)
 \quad(x\ge2).                                         \tag{NT-A04a}
\]
Any endpoint convention is allowed. In particular, uniformly for
$z\in\mathbb C$,
\[
 \#\{\mathfrak p\text{ unramified}:N\mathfrak p\le x,\
                                  \eta(\mathfrak p)=z\}
       \ll_{K,\eta}\sqrt x\log^2(2x).                    \tag{NT-A04b}
\]
The hypothesis on the character is essential for the angular law;
infinite order by itself is insufficient.

\paragraph{Exact classical locator and deductions.}
Rajan, \emph{Distribution of values of Hecke characters of infinite
order}, Acta Arithmetica85 (1998), Proposition4, printed
pp.289--290, gives (NT-A04a) with $\pi_K(x)$ in place of
$\operatorname{Li}_2(x)$. The character restriction is the one fixed
in Theorem1, printed p.286: in that paper a \emph{Dirichlet character}
means a finite-order Hecke character of $K$, and equivalence means
equality up to a power of the norm. The arguments use normalized
angles modulo one. Proposition4 has a constant depending only on
$K$ above a character-dependent threshold; allowing the constant
to depend on the fixed $\eta$ gives the displayed range $x\ge2$.
The finite ramification set is omitted throughout. Its analytic
premises are ordinary Hecke continuation and the precisely stated
power RH family; $M$ supplies them conditionally.

The prime-ideal estimate NT-C02 for the trivial extension $K/K$
replaces $\pi_K(x)$ by $\operatorname{Li}_2(x)$ with an error
$O_K(\sqrt x\log(2x))$. Shrinking or enlarging an arc and passing
to the limit for the finite collection of primes handles all
endpoint conventions, exactly as in NT-E06. For $|z|=1$, shrink
arcs to the singleton $z$; their main terms tend to zero while the
uniform error is fixed. This proves (NT-A04b). For $|z|\ne1$ the
set is empty.

\paragraph{Complete check of the norm-twist obstruction.}
Suppose instead that $\eta(\mathfrak p)=
\chi(\mathfrak p)(N\mathfrak p)^{it}$, where $\chi$ has finite order
$h$ and $t\ne0$. Set $\tau=ht\ne0$. Then the $h$th Fourier moment
of its prime angles is
\[
 \frac1{\pi_K(x)}\sum_{N\mathfrak p\le x}\eta(\mathfrak p)^h
       =\frac{x^{i\tau}}{1+i\tau}+o_{K,\chi,t}(1).       \tag{NT-A04c}
\]
Omitting the fixed ramification set does not affect this formula.
Here is a direct proof under the same available input $M$.
By NT-C02,
$\pi_K(u)=\operatorname{Li}_2(u)+O_K(\sqrt u\log(2u))$.
Abel summation therefore gives
\[
 \sum_{N\mathfrak p\le x}(N\mathfrak p)^{i\tau}
 =\int_2^x\frac{u^{i\tau}}{\log u}\,du
      +O_{K,\tau}(\sqrt x\log(2x)+1).
\]
Integration by parts with primitive
$u^{1+i\tau}/(1+i\tau)$ shows that the integral is
\[
 \frac{x^{1+i\tau}}{(1+i\tau)\log x}
           +O_\tau\!\left(\frac{x}{\log^2 x}+1\right).
\]
For the remainder, bound
$\int_2^xdu/\log^2u\ll x/\log^2x$, splitting at $\sqrt x$ if
desired. Since $\pi_K(x)\sim x/\log x$, (NT-A04c) follows.
Its absolute value tends to $(1+\tau^2)^{-1/2}>0$, whereas a
uniform angular distribution would force every nonzero Fourier
moment to tend to zero. Thus even under $M$ these infinite-order
characters do not obey (NT-A04a). The source exclusion prevents
the norm-twist poles from being silently treated as entire
nonprincipal factors.

\paragraph{Geometric dictionary and location.}
The general statement concerns rank-one unitary characters of the
idele class group. If the character is the unitary normalization
of an algebraic Hecke character, this is the Frobenius-value
distribution of the associated rank-one CM motive and its tensor
powers. A general unitary infinity type is not automatically
asserted to be an algebraic or \'{e}tale motive.
Work/run: \texttt{GRH\_NIGHT/number\_theory},
\texttt{hecke\_shrinking\_angles.tex}, NT-A04; primary source
S-NT-RAJAN, with the complete exception proof above.

\subsection*{NT-A05. Normalized shrinking arcs for infinite-order Hecke values}
Under the exact assumptions of NT-A04, fix
$c>0$ and $0<\alpha<1/2$, and set
$u_0=\max\{2,c^{1/\alpha}\}$.
For a center $\vartheta\in\mathbb R/\mathbb Z$ and $u\ge u_0$,
let $J_u(\vartheta)$ be the arc centered at $\vartheta$ of length
$cu^{-\alpha}$, with any fixed endpoint convention. Then, as $x\to\infty$,
\begin{align}
 &\#\{\mathfrak p\text{ unramified}:u_0<N\mathfrak p\le x,\
               \theta_{\mathfrak p}\in J_{N\mathfrak p}(\vartheta)\}
       \notag\\
 &\hspace{12mm}
   =c\int_{u_0}^x\frac{u^{-\alpha}}{\log u}\,du
       +O_{K,\eta,c,\alpha}
          \bigl(x^{3/4-\alpha/2}\log^2(2x)\bigr)
       \notag\\
 &\hspace{12mm}\sim\frac{c}{1-\alpha}\frac{x^{1-\alpha}}{\log x}.
                                                               \tag{NT-A05}
\end{align}
The error constant is uniform in the chosen center $\vartheta$.
All other parameters are fixed, and no assertion at $\alpha=1/2$
is made.

\paragraph{Complete proof.}
Set $h=x^{-1/4+\alpha/2}$. Partition $[u_0,x]$ into
successive multiplicative intervals of ratio $1+h$, with the
last endpoint truncated at $u_0$. The arcs are nested for the
fixed center. On each norm interval the moving prime count is
bounded above and below by the counts for its two endpoint arcs.
Apply (NT-A04a) at both norm endpoints, with a common constant
independent of $\vartheta$. The sum of all endpoint errors is
\[
 \ll_{K,\eta}\log^2(2x)
      \sum_{j\ge0}\sqrt{x(1+h)^{-j}}
 \ll_{K,\eta}\sqrt x\,h^{-1}\log^2(2x).
\]
The main brackets are the upper and lower Riemann sums for
$c\int_{u_0}^x u^{-\alpha}\,du/\log u$.
On an interval of ratio at most $1+h$, the endpoint values of
$u^{-\alpha}$ differ by a factor at most
$(1+h)^\alpha=1+O_\alpha(h)$. The difference of these Riemann
sums is consequently
$O_{c,\alpha}(h x^{1-\alpha})$.
With the chosen $h$, both power terms are $x^{3/4-\alpha/2}$.
This proves the equality in (NT-A05), with uniformity in the
center and the stated endpoint conventions.

Finally, integration by parts gives
$\int_{u_0}^x u^{-\alpha}\,du/\log u
\sim x^{1-\alpha}/((1-\alpha)\log x)$.
The ratio of the error to that main term is
$O(x^{-1/4+\alpha/2}\log^3(2x))=o(1)$.
This proves the asymptotic and the exact open range
$0<\alpha<1/2$. The window is normalized in the compact rank-one
torus; it is an arithmetic prime-counting assertion, and it does
not identify zeros of the Hecke function with independent random
angles. Work/run: \texttt{GRH\_NIGHT/number\_theory},
\texttt{hecke\_shrinking\_angles.tex}, NT-A05.
